\documentclass[aspectratio=169,11pt]{beamer}
\usepackage[T1]{fontenc}
\usepackage[utf8]{inputenc}
\usepackage[ngerman]{babel}
\usepackage{lmodern,amsmath,amssymb,booktabs,array,tabularx}
\usepackage{pgfplots}
\pgfplotsset{compat=1.18}
\definecolor{navy}{HTML}{17324D}
\definecolor{teal}{HTML}{007F82}
\definecolor{rust}{HTML}{BC623B}
\definecolor{muted}{HTML}{586777}
\setbeamercolor{normal text}{fg=navy,bg=white}
\setbeamercolor{frametitle}{fg=navy}
\setbeamercolor{structure}{fg=teal}
\setbeamercolor{alerted text}{fg=rust}
\setbeamerfont{frametitle}{size=\Large,series=\bfseries}
\setbeamerfont{normal text}{size=\normalsize}
\setbeamersize{text margin left=8mm,text margin right=8mm}
\setbeamertemplate{navigation symbols}{}
\setbeamertemplate{frametitle}{\vspace{3mm}\insertframetitle\par\vspace{1mm}}
\setbeamertemplate{footline}{\hspace{8mm}\color{muted}\scriptsize freMTPL2\hfill\insertframenumber\hspace{8mm}\vspace{2mm}}
\setbeamertemplate{itemize items}[circle]
\hypersetup{colorlinks=true,urlcolor=teal,linkcolor=teal,pdftitle={Von der Exponierung zum erwarteten Schadenaufwand},pdfauthor={freMTPL2-Portfolio}}
\newcommand{\source}[1]{\vfill{\fontsize{6.5}{7.4}\selectfont\color{muted}#1\par}}
\newcommand{\smalltext}{\fontsize{9}{11.3}\selectfont}
\newcommand{\takeaway}[1]{\par\medskip{\color{teal}\bfseries #1}\par}
\newcolumntype{Y}{>{\raggedright\arraybackslash}X}
\newcolumntype{L}[1]{>{\raggedright\arraybackslash}p{#1}}
% Generated from aggregate reports; no model calls.
\newcommand{\EvalGain}{9{,}80}
\newcommand{\EvalObserved}{173{,}58}
\newcommand{\EvalPredicted}{165{,}60}
\newcommand{\EvalAP}{1{,}048}
\newcommand{\EvalPolicies}{135\,599}
\newcommand{\ReducedGap}{3{,}5773}
\newcommand{\AgeBMRho}{-0{,}444}
\newcommand{\AgeBMRank}{-0{,}572}
\newcommand{\AreaAgeEta}{0{,}00260}
\newcommand{\AreaBMEta}{0{,}01672}
\newcommand{\LargestShare}{8{,}58}
\newcommand{\LargestClaim}{4\,075\,400{,}56}
\newcommand{\AugMin}{0{,}37}
\newcommand{\AugMax}{3{,}06}
\newcommand{\AugBaseline}{0{,}464386}
\newcommand{\SupportMin}{73{,}77}
\newcommand{\SupportMax}{80{,}40}
\newcommand{\FrequencyTable}{%
\begin{tabular}{lrr}\toprule
Regel & $D_N$ & Gewinn\,\%\\\midrule
Intercept & 0{,}478094 & 0{,}00\\
GLM & 0{,}457123 & 4{,}39\\
GAM & 0{,}454637 & 4{,}91\\
HGB & 0{,}446749 & 6{,}56\\
\bottomrule\end{tabular}
}
\newcommand{\LossTable}{%
\begin{tabular}{lrr}\toprule
Schadenhöhe & OOF $D_L$ & Eval. $D_L$\\\midrule
Mittel (gewählt) & 79{,}195 & 80{,}207\\
Gamma-GLM & 80{,}387 & 80{,}064\\
Log-Smearing & 79{,}370 & 80{,}026\\
\bottomrule\end{tabular}
}
\newcommand{\FeatureTable}{%
\begin{tabular}{lr}\toprule
Regel & OOF $D_L$\\\midrule
Alle Merkmale (HGB) & 79{,}195\\
Alter + BM (GAM) & 82{,}028\\
BM allein (GAM) & 82{,}537\\
BM + Area (GAM) & 82{,}621\\
\bottomrule\end{tabular}
}
\newcommand{\AllFeatureTable}{%
\begin{tabular}{lrr}\toprule
Merkmale & GLM $D_L$ & GAM $D_L$\\\midrule
Alter & 84{,}717 & 84{,}318\\
BM & 83{,}240 & 82{,}537\\
Area & 87{,}283 & 87{,}280\\
BM + Area & 83{,}389 & 82{,}621\\
Alter + Area & 84{,}620 & 84{,}181\\
Alter + log(Density) & 84{,}504 & 84{,}057\\
Alter + Region & 84{,}668 & 84{,}216\\
Alter + BM & 82{,}621 & 82{,}028\\
\bottomrule\end{tabular}
}
\newcommand{\AgeProfileTable}{%
\begin{tabular}{lrrr}\toprule
Refit & 35 Jahre & 45 Jahre & 55 Jahre\\\midrule
GLM & 0{,}06451 & 0{,}08177 & 0{,}07227\\
GAM & 0{,}05249 & 0{,}07773 & 0{,}07411\\
HGB & 0{,}05588 & 0{,}07707 & 0{,}07233\\
\bottomrule\end{tabular}
}
\newcommand{\GeoTable}{%
\begin{tabular}{lrr}\toprule
Area & Rate/Jahr & EUR/Schaden\\\midrule
AB & 0{,}0572 & 2\,879\\
CD & 0{,}0743 & 2\,079\\
EF & 0{,}0948 & 2\,050\\
\bottomrule\end{tabular}
}
\newcommand{\StandardTable}{%
\begin{tabular}{llrr}\toprule
Mix & Area & Rate/Jahr & EUR/Jahr\\\midrule
ohne Dauer & AB & 0{,}05958 & 177{,}15\\
ohne Dauer & CD & 0{,}07176 & 149{,}13\\
ohne Dauer & EF & 0{,}08647 & 177{,}83\\
mit Dauer & AB & 0{,}05643 & 112{,}51\\
mit Dauer & CD & 0{,}06732 & 137{,}59\\
mit Dauer & EF & 0{,}08051 & 154{,}70\\
\bottomrule\end{tabular}
}
\newcommand{\FrequencyDuration}{%
\begin{tabular}{lr}\toprule
Vertragsdauer & A/P\\\midrule
7 Tage--<0,1 J. & 1{,}794\\
0,1--<0,5 J. & 1{,}237\\
0,5--<1 J. & 0{,}998\\
1+ J. & 0{,}833\\
\bottomrule\end{tabular}
}
\newcommand{\LossDuration}{%
\begin{tabular}{lr}\toprule
Vertragsdauer & A/P\\\midrule
7 Tage--<0,1 J. & 3{,}057\\
0,1--<0,5 J. & 2{,}152\\
0,5--<1 J. & 0{,}797\\
1+ J. & 0{,}607\\
\bottomrule\end{tabular}
}
\newcommand{\AgeLabels}{18--24,25--29,30--34,35--39,40--44,45--49,50--54,55--59,60--64,65--74,75+}

\pgfplotsset{every axis/.append style={axis line style={muted!50},tick style={muted!50},
 tick label style={font=\scriptsize},label style={font=\scriptsize},
 title style={font=\small\bfseries},grid=major,grid style={gray!15},
 axis lines=left,scaled ticks=false,/pgf/number format/use comma,/pgf/number format/fixed}}

\begin{document}
\ifdefined\BackupOnly\else

% 1 / 0:50
\begin{frame}[t]{Von der Exponierung zum Schadenaufwand}
\vspace{3mm}
{\large Wie viel Modellkomplexität hilft bei freMTPL2?}\par
\medskip
\[
\underbrace{q(x)}_{\text{EUR/Vertragsjahr}}
=\underbrace{r(x)}_{\text{Schäden/Vertragsjahr}}
\;\times\;
\underbrace{m(x)}_{\text{EUR/Schaden}}
\]
\medskip
\textbf{HGB-Frequenz mit konstantem Schadenmittel:}
\EvalGain\,\% weniger Verlustdeviance als die Interceptbasis
in der bestehenden Evaluation.
\medskip

Das konstante Schadenmittel erzielt den kleinsten Entwicklungs-Verlustscore.
Im CTGAN-Pilot schneiden beide Kontrollen besser ab.
\source{Kernstudie: reports/s6/results.json; CTGAN-Pilot: reports/o2/report.json.\quad
Historischer französischer Bestand; kein verwendungsfertiger Markttarif.\quad Fassung vom 9. November 2026}
\end{frame}

% 2 / 1:00
\begin{frame}[t]{Ein abgeglichener französischer Haftpflichtbestand}
Unbekannter privater Versicherer, ungefähr 2011--2013.\par
\medskip
\begin{center}
\begin{tabular}{lrr}\toprule
& CASdatasets 1.2-1 & OpenML v1\\\midrule
Vertragszeilen & 677\,991 & 678\,013\\
Mitgelieferte Schadenanzahl & 26\,444 & 36\,102\\
Verknüpfte positive Schadenzeilen & 26\,444 & 26\,444\\
Nicht zuordenbare Schadenzeilen & 0 & 195\\\bottomrule
\end{tabular}
\end{center}
\medskip
\textbf{Primärziel:} erfasste positive Schadenzeilen und zugehörige
historische Zahlbeträge; Nullschadenverträge bleiben enthalten.
\smallskip

CAS: 358\,482,84 Vertragsjahre und 59,91 Mio. EUR.
Anzahl, Exponierung und Beträge bleiben ungekappt.
\source{reports/s0/source\_comparison.json; CASdatasets-Commit 227fb56b8734, Dokumentation [Q1].
Bei 9\,095 gemeinsamen Verträgen weicht die mitgelieferte Anzahl ab; die Ursache ist ungeklärt.
Area/Density: Wohnort; Region: Vertragsregion. Teilweise IRSA-IDA-Konventionsbeträge.}
\end{frame}

% 3 / 1:25
\begin{frame}[t]{Der Altersverlauf hängt auch vom Bestandsmix ab}
\begin{columns}[T,onlytextwidth]
\column{.58\textwidth}
\begin{tikzpicture}
\begin{axis}[width=8.0cm,height=3.8cm,title={Beobachtete Frequenz},
 ymin=0.04,ymax=0.18,xmin=.5,xmax=11.5,xtick={1,...,11},
 xticklabels={18--24,25--29,30--34,35--39,40--44,45--49,50--54,55--59,60--64,65--74,75+},
 x tick label style={rotate=55,anchor=east,font=\tiny},ylabel={Schäden/Vertragsjahr}]
\addplot+[teal,only marks,mark=*,error bars/.cd,y dir=both,y explicit]
 table[x=x,y=rate,y error minus=low,y error plus=high]{age.dat};
\end{axis}
\end{tikzpicture}
\column{.40\textwidth}
\begin{tikzpicture}
\begin{axis}[width=5.4cm,height=3.8cm,title={Gemeinsame BM-Mischung},
 ymin=.04,ymax=.10,xmin=.7,xmax=3.3,xtick={1,2,3},
 xticklabels={30--39,40--49,50--59},
 legend style={at={(.5,-.3)},anchor=north,font=\tiny,draw=none}]
\addplot+[muted,mark=square*] table[x=x,y=raw]{bm.dat};\addlegendentry{marginal, gleicher Support}
\addplot+[teal,mark=*,error bars/.cd,y dir=both,y explicit]
 table[x=x,y=standard,y error minus=low,y error plus=high]{bm.dat};\addlegendentry{BM-standardisiert, 95\,\%}
\end{axis}
\end{tikzpicture}
\end{columns}
\vspace{3mm}
\smalltext Links: Entwicklung, 542\,392 Verträge. Rechts: nur gemeinsame BM-Unterstützung;
andere Population. Der Wiederanstieg verschwindet dort nicht.
\source{reports/s3/age\_bootstrap.json, bm\_standardized.json (gepinnt); punktweise Intervalle aus 500 Vertragsziehungen, kein simultanes Band.
Die mittleren Fünfjahresgruppen tragen jeweils rund 35\,000 Vertragsjahre; lokale Nachbarkontraste sind unsicher.
Support und weitere Adjustierung: B6.}
\end{frame}

% 4 / 1:30
\begin{frame}[t]{Geografie: mehr Frequenz, unsichere Schadenhöhe}
\begin{columns}[T,onlytextwidth]
\column{.43\textwidth}
\textbf{freMTPL2, Entwicklung}\par\medskip
{\small\GeoTable}\par\smallskip
\smalltext AB: ländlicher, EF: urbaner.\par
EF--AB nach Standardisierung mit Vertragsdauer:
0,0241 [0,0210; 0,0272] Schäden/Jahr.
\smallskip

Beim Unterschied der Schadenmittel (EF--AB) schließen die Intervalle null ein.
\column{.54\textwidth}
\textbf{Was externe Evidenz beiträgt}\par\medskip
\smalltext
\textbf{ONISR 2024 [Q2]:} Auf Außerortsstraßen ohne Autobahnen entfallen 60\,\% der Toten und 47\,\% der geschätzten Schwerverletzten.
\par\smallskip
\textbf{Cerema [Q3]:} Nach der 80-km/h-Regel sinken Geschwindigkeit und Mortalität im Straßenvergleich.
Das stützt Geschwindigkeit als Mechanismus körperlicher Schwere.\par\smallskip
Unfallort und Verletzungsschwere unterscheiden sich von Wohnort und versichertem Zahlbetrag.
\end{columns}
\takeaway{Der Frequenzkontrast bleibt; seine Ursache ist nicht identifiziert}
\source{reports/s6a/geo.json, standardized\_with\_duration.json und zugehörige Intervalle; bedingte 95\,\%-Intervalle.
Auch adjustierte Altersprofile bleiben nicht monoton (B6).
Originalquellen, Vergleichsdesign und Populationen: B8, Q1--Q3.}
\end{frame}

% 5 / 1:00
\begin{frame}[t]{Drei Modellklassen für denselben Frequenzmittelwert}
\[
\mu_i=e_i\,r(x_i),\qquad \log\mu_i=\underbrace{\log e_i}_{\text{Offset}}+\eta(x_i)
\]
\medskip
\begin{tabularx}{\textwidth}{L{20mm}L{58mm}Y}\toprule
Modell & Struktur & Welche Frage beantwortet es?\\\midrule
GLM & Gruppen, Faktoren, log-Dichte; auch Alter--BM-Interaktion & Reichen einfache Relativitäten?\\[5pt]
GAM & $\eta=\beta_0+z^\top\gamma+\sum_j B_j(x_j)\theta_j$; zweite Differenzenstrafe & Helfen glatte nichtlineare Verläufe?\\[5pt]
HGB & Poisson-Boosting, bis zu 15 oder 31 Endknoten je Baum & Was fehlt der additiven Struktur?\\\bottomrule
\end{tabularx}
\medskip
Alle lernen $N_i/e_i$ mit Gewicht $e_i$; Vorverarbeitung nur im Trainingsfold.
\source{config/frequency\_protocol.json; src/frequency\_models.py; [M1--M3].
300 Boosting-Iterationen. Das GAM mit zwei Merkmalen nutzt eine andere Strafe (B3).}
\end{frame}

% 6 / 1:15
\begin{frame}[t]{Boosting senkt die Deviance. Die Dauer bleibt wichtig.}
\smalltext 80\,\% Entwicklung mit fünf OOF-Folds; 20\,\% gemeinsame Evaluation erst nach Auswahl.
Historisch bereits explorierte Daten: keine externe oder zeitliche Prüfung.\par\medskip
\begin{columns}[T,onlytextwidth]
\column{.50\textwidth}
\textbf{Entwicklungs-OOF}\par\medskip
{\small\FrequencyTable}\par\medskip
\smalltext GAM minus HGB: 0,00789\par
95\,\%-Intervall [0,00618; 0,00965].\par
GLM enthält im Familienminimum die Alter--BM-Interaktion.
\column{.47\textwidth}
\textbf{HGB-Frequenz-A/P (Entwicklungs-OOF)}\par\medskip
{\small\FrequencyDuration}\par\medskip
\smalltext A/P = tatsächliche / erwartete Anzahl.\par
Gesamt fast eins; deutliche Dauerabweichungen trotz unterstützter Gruppen.
\end{columns}
\source{reports/s6/results.json; reports/s4/scores.json, paired\_bootstrap.json, segments.json.
$D_N$: Poisson-Deviance je Vertragsjahr, kleiner ist besser; Gewinn relativ zum Intercept.
Intervalle: 500 gepaarte Vertragsziehungen bei festen OOF-Vorhersagen, ohne Refit-/Auswahlunsicherheit.}
\end{frame}

% 7 / 1:15
\begin{frame}[t]{Zusätzliche Schadenhöhenmerkmale helfen hier nicht sicher}
\begin{columns}[T,onlytextwidth]
\column{.56\textwidth}
\textbf{HGB-Frequenz $\times$ Schadenhöhe}\par\medskip
{\small\LossTable}\par\medskip
\smalltext $D_L$: Jahresraten-Tweedie-Deviance, $p=1{,}5$.\par
Die Wahl fällt in der Entwicklung auf das Mittel.
Kleine umgekehrte Evaluationsabstände rechtfertigen keine Neuwahl;
ihre gepaarten Intervalle enthalten null.
\column{.40\textwidth}
\textbf{Gewähltes Produkt, Evaluation}\par\medskip
\EvalPredicted\ EUR/Jahr prognostiziert\par
\EvalObserved\ EUR/Jahr beobachtet\par
\medskip
Verlust-A/P: \EvalAP\par\medskip
\smalltext Entwicklung: Ein einziger Schaden trägt \LargestShare\,\% der Verlustsumme.
\end{columns}
\medskip
\[
\widehat Q=\frac{\sum_i e_i\widehat r_i\widehat m_i}{\sum_i e_i}
\qquad\text{erst Vertragsprodukte bilden, dann aggregieren}
\]
\source{reports/s6/results.json; reports/s5/evaluation\_bootstrap.json; reports/s3/tail\_concentration.json (gepinnt).
Evaluation: \EvalPolicies\ Verträge. Gamma verbessert dort den separaten Claim-Score;
kein allgemeines Urteil gegen Schadenhöhenmodelle. Produktannahme: B5.
Verlust-A/P: 3,057 bei 7 Tagen--unter 0,1 Jahren, 0,607 ab einem Jahr.}
\end{frame}

% 8 / 1:30
\begin{frame}[t]{Alter und Bonus-Malus: Prognosewert trotz Überlappung}
\begin{columns}[T,onlytextwidth]
\column{.51\textwidth}
{\small\FeatureTable}\par\medskip
\smalltext Alter + BM: rund 3,58\,\% höhere Deviance als das vollständige Modell.
BM + Area minus Alter + BM:\par
0,593 [--0,144; 1,803].\par $D_L$: exponierungsgewichtete Tweedie-Deviance der Jahresraten.
\column{.46\textwidth}
\smalltext
\textbf{Zusatznutzen von Alter}\par
Alter neben BM verbessert den Verlustscore um 0,509 [0,115; 1,103].\par\medskip
\textbf{Überlappung von Alter und BM}\par
Exponierungsgewichtete Alters-BM-Korrelation: \AgeBMRho.\par\medskip
\textbf{Stabilität der Einzeleffekte}\par
Intervalle bei festen Fits erfassen keine Streuung bei erneuter Anpassung. Auch für Area ist die Stabilität offen.
\end{columns}
\takeaway{Prognoseleistung, Erklärbarkeit und stabile Einzeleffekte gemeinsam abwägen}
\source{reports/s6a/scores.json und ablation/\{scores,dependence,bootstrap\_two\_BonusMalus\_gam\}.json.
Explorativer Vergleich auf den Entwicklungs-OOF-Folds.
500 Vertragsziehungen bei festen Fits; keine globale Optimalität des Merkmalsduos. Vollständiger Vergleich: B9.}
\end{frame}

% 9 / 1:05
\begin{frame}[t]{Synthetische Trainingszeilen liefern im Pilot keinen Nutzen}
\begin{columns}[T,onlytextwidth]
\column{.56\textwidth}
\begin{tikzpicture}
\begin{axis}[width=7.4cm,height=4.7cm,xmin=.5,xmax=6.5,ymin=0,ymax=.018,
 ytick={0,.005,.010,.015},y tick label style={font=\scriptsize,/pgf/number format/precision=3,/pgf/number format/fixed zerofill},
 ylabel={$\Delta D_N$ je Vertragsjahr},xtick={1,...,6},
 xticklabels={{0,25/17},{0,25/29},{0,25/43},{1/17},{1/29},{1/43}},
 x tick label style={font=\tiny,rotate=30,anchor=east},xlabel={Zusatzratio / Seed},
 legend style={at={(.02,.98)},anchor=north west,font=\tiny,draw=none}]
\addplot+[teal,only marks,mark=*,error bars/.cd,y dir=both,y explicit]
 table[x=x,y=real,y error minus=low,y error plus=high]{augmentation.dat};\addlegendentry{CTGAN minus real-only, 95\,\%}
\addplot+[rust,only marks,mark=square*] table[x=x,y=bootstrap]{augmentation.dat};\addlegendentry{CTGAN minus Bootstrap, Punktwert}
\end{axis}
\end{tikzpicture}
\column{.41\textwidth}
\smalltext Fünf bestehende Folds, je 20\,000 reale Trainingsverträge;
drei Seeds, 30 Epochen. Zusätzliche Zeilen im Verhältnis 0,25 oder 1.\par\medskip
Alle sechs CTGAN-Varianten sind schlechter als beide Kontrollen:
\textbf{+\AugMin--\AugMax\,\%} gegenüber real-only.\par\medskip
Begrenztes Budget, fester GLM. Kein Test der Augmentation des vollständigen Bestands und kein allgemeines CTGAN-Urteil.
\end{columns}
\source{reports/o2/report.json; reale Entwicklungs-OOF-Verträge, Basis $D_N=\AugBaseline$.
Intervalle: gepaarte asymptotische 95\,\%-Intervalle bei festen Vorhersagen; keine Generator-/Refitunsicherheit.
Sampling, gemeinsame Unterstützung und Datenschutzgrenzen: B11; Xu et al. [M4].}
\end{frame}

% 10 / 0:50
\begin{frame}[t]{Was die Studie trägt und was vor Tarifverwendung fehlt}
\vspace{2mm}
\begin{tabularx}{\textwidth}{L{55mm}Y}\toprule
Tragfähiger Befund & Nächster notwendiger Nachweis\\\midrule
Flexible Frequenz verbessert den historischen Score & Zeitliche und externe Prüfung auf einem aktuellen Zielbestand\\[8pt]
Wenige Merkmale liefern viel Prognoseinformation & Refit-/Effektstabilität, Segmentkalibrierung und gemeinsame Unterstützung\\[8pt]
Kurze Laufzeiten sind schlecht kalibriert. Großschäden prägen das Schadenmittel. & Fahrleistung, Nutzung, Beginn/Ende, Schaden- und Regulierungsverlauf\\\bottomrule
\end{tabularx}
\takeaway{Ein nachvollziehbares Schadenbedarfsmodell ist die Grundlage; der Marktpreis braucht mehr Evidenz}
\source{reports/s6/results.json; docs/aktuarische\_einordnung.md; reports/s6a/ablation/report.json.
Für Deutschland zusätzlich: Deckung, Haftung, Inflation und Zielmix; französischer Bonus-Malus ist keine deutsche SF-Klasse.
Kosten, Vertrieb und Transfer: B12.}
\end{frame}

\fi
\ifdefined\MainOnly\else
\appendix
\setcounter{framenumber}{0}
\setbeamertemplate{footline}{\hspace{8mm}\color{muted}\scriptsize freMTPL2 | Backup\hfill B\insertframenumber\hspace{8mm}\vspace{2mm}}

\begin{frame}[t]{B1 | Offset, Rate und Periodenerwartung}
\[
N_i\sim\operatorname{Poisson}(e_i\exp\eta_i),\quad
\mu_i=e_i r_i,\quad \log\mu_i=\log e_i+\eta_i
\]
\[
-\ell=\sum_i(e_i\exp\eta_i-N_i\eta_i)+C(N,e)
=\sum_i e_i\{\exp\eta_i-(N_i/e_i)\eta_i\}+C
\]
\textbf{Rate als Ziel, Exponierung als Gewicht:} dieselben parameterabhängigen Terme.
Gebrochene Raten sind keine behaupteten Poisson-Anzahlen.
\medskip

$r_i$: Schäden/Jahr; $\mu_i$: erwartete Anzahl der tatsächlichen Periode.
Eine Jahresrate ist keine Schadenwahrscheinlichkeit.
\medskip

Unbestrafter Intercept: $\sum_iN_i=\sum_i e_i\widehat r_i$ am geeigneten Optimum.
Trainingsbalance garantiert keine OOF- oder Segmentkalibrierung.
\source{paper/begleitpaper.tex, Abschnitte 3--4; src/frequency\_models.py; [M1].
Arbeitsverteilung versus Mittelwertmodell unterscheiden; Unabhängigkeit für eine Produktlikelihood ist eine zusätzliche Annahme.}
\end{frame}

\begin{frame}[t]{B2 | Welche Verlustmaße werden verglichen?}
\[
d(N,\mu)=2\{N\log(N/\mu)-N+\mu\},\quad 0\log0=0
\]
\[
D_N=\frac{\sum_i d(N_i,e_i\widehat r_i)}{\sum_i e_i}
=\frac{\sum_i e_i d(N_i/e_i,\widehat r_i)}{\sum_i e_i}
\]
\[
D_L=\frac{\sum_i e_i d_{1.5}(S_i/e_i,\widehat q_i)}{\sum_i e_i},\qquad
d_{1.5}(y,q)=4\frac{(\sqrt y-\sqrt q)^2}{\sqrt q}
\]
\smalltext Bei Tweedie gilt $d_p(e y,e q)=e^{2-p}d_p(y,q)$.
Eine mittlere Periodendeviance gewichtet daher anders als $D_L$.
Scores sind weder EUR-Fehler noch Prozentsätze; nur ihre relative Änderung wird in Prozent berichtet.\par\medskip
Poisson-Scoredifferenzen sind bei endlichem positiven bedingten Mittel mittelwertkonsistent.
Für einzelne erwartete Deviances genügt zusätzlich $E[N\log_+N\mid x,e]<\infty$.
Bei Mittel null liegt das Optimum am Rand.
\source{src/metrics.py, src/severity.py; paper Abschnitte 3--4 und 8; Gneiting (2011), [M5].
Für Kalibrierung gilt A/P = Summe beobachtet / Summe prognostiziert; keine zusätzliche Exponierungsgewichtung von Anzahlen.}
\end{frame}

\begin{frame}[t]{B3 | Strafen, Identifizierung und effektive Freiheitsgrade}
\smalltext
\textbf{Vollständiges GLM:}
$\sum_i(e_i\exp\eta_i-N_i\eta_i)/\sum_i e_i+\alpha\|b\|^2/2$;
Intercept unbestraft. Die Summe der Gewichte normiert den Verlust.\par\medskip
\textbf{GAM im vollständigen Frequenzvergleich:}
Poisson-Summe $+\lambda\sum_j\|D_2\theta_j\|^2/2$.
Kubische B-Splines, 10/20 Basen, $\lambda=1/10$;
Faktoren unbestraft. Numerische Strafparameter sind nicht mit GLM-$\alpha$ vergleichbar.\par\medskip
\textbf{Identifizierung:} letzte Basisfunktion entfernen, übrige mit Trainings-Exponierung zentrieren;
Strafe in derselben Parametrisierung abbilden. $\sum_i e_i f_j(x_{ij})=0$.
\[
\operatorname{EDF}=\operatorname{tr}\{(H+P)^{-1}H\}
\]
EDF gelten bedingt auf die gewählte Glättung; Rang des unbestraften Designs
und positive Definitheit von $H+P$ sind unterschiedliche Prüfungen.\par\medskip
\textbf{GAM mit zwei Merkmalen:} acht uniforme Knoten, kubisch und zentriert;
Ridge $\alpha=10^{-4}$ auf nichtkonstante Koeffizienten, keine $D_2$-Strafe.
\source{config/frequency\_protocol.json, interpretation\_protocol.json; src/frequency\_models.py, src/interpretation.py.
Eilers \& Marx (1996), Wood, [M2--M3]. Randfortsetzung konstant; isolierte Effekte bei korrelierten Merkmalen vorsichtig lesen.}
\end{frame}

\begin{frame}[t]{B4 | Schadenhöhe: Gewicht eins und echtes log(Y)-Smearing}
\textbf{Population:} positive Trainingsschadenzeilen; jede Zeile Gewicht eins.
Alle Schäden eines Vertrags bleiben in derselben Aufteilung.
Kein Exponierungsoffset und keine Exponierungsgewichte.
\medskip

\smalltext Trainingsclaim-Mittel: $\widehat m=\sum_kY_k/n_{\rm Claims}$.
Gamma-GLM mit Log-Link: Verlust proportional zu $\sum_k(Y_k/\exp\eta_k+\eta_k)$.
Seine Interceptgleichung erzwingt $\operatorname{mean}(Y/\widehat m)=1$, keine Euro-Summenbalance.\par\medskip
Implementiert ist $\log Y=\eta(x)+\varepsilon$, nicht $\log(1+Y)$:
\[
\widehat m(x)=\exp(\widehat\eta(x))\,
\underbrace{\frac1{n_{\rm train}}\sum_{k\in\rm train}
 \exp(\log Y_k-\widehat\eta(x_k))}_{\text{globaler Trainings-Smearingfaktor}}
\]
Kein Abzug von eins. Bedingt korrekt wäre der Faktor
$E[\exp\varepsilon\mid x]$; seine globale Ersetzung ist eine Annahme.
Nullmittelwert von Log-Residuen reicht dafür nicht.
\source{src/severity\_models.py; config/severity\_protocol.json; paper Abschnitt 7; Duan (1983), [M6].
Dieselbe prospektive Funktion prognostiziert alle Verträge, auch spätere Nullschäden.
Claim-Gamma-Deviance und Verlustscore beurteilen unterschiedliche Populationen.}
\end{frame}

\begin{frame}[t]{B5 | Produktannahme, Vertragssummen und Tail}
\[
E[S\mid N,x,e]=N m(x),\quad E[N\mid x,e]=e r(x)
\ \Longrightarrow\ E[S\mid x,e]=e r(x)m(x)
\]
\smalltext Hinreichende bedingte Mittelwertannahme und endliche beteiligte Mittel;
der Datenabgleich beweist sie nicht. Vollständige Unabhängigkeit aller Beträge ist nicht nötig.\par\medskip
Erwartete Eurosummen: $\sum_i e_i\widehat r_i\widehat m_i$.
Jahresbedarf: diese Summe geteilt durch $\sum_i e_i$.
Getrennte exponierungsgewichtete Komponentenmittel ergeben im Allgemeinen ein anderes Produkt.
Auch konstante Schadenmittel können zwischen OOF-Folds variieren.\par\medskip
\textbf{Tail:} größter Entwicklungsschaden \LargestClaim\ EUR,
\LargestShare\,\% der Verlustsumme.
Kappung ist ein anderes Ziel:
\[
E[Y]-E[\min(Y,u)]=E[(Y-u)_+]=\int_u^\infty\Pr(Y>t)\,dt
\]
\source{paper Abschnitt 8; reports/s3/tail\_concentration.json (gepinnt).
Ungekappte Hauptanalyse; empirischer Bootstrap kann unbekannte Extremverluste nicht erzeugen.}
\end{frame}

\begin{frame}[t]{B6 | Altersprofile, Unterstützung und Beobachtung}
\begin{center}\small\AgeProfileTable\end{center}
\smalltext Jahresraten bei gesetztem Alter 35/45/55 auf identischen Referenzverträgen;
Profile aus Fits auf der vollständigen Entwicklung.\par\smallskip
291\,234 Verträge, 163\,403,34 Vertragsjahre; Alter 30--59.
Gemeinsame BM-/Fahrzeugalter-/Kraftstoff-/Regionszellen mit mindestens 100 Vertragsjahren je Altersdekade.
Kein Nachweis exakten gemeinsamen Supports.\par\smallskip
Keine rein monotone Erfahrungskurve; Erfahrung, Nutzung und Unfallfahrer sind nicht getrennt beobachtet.
Drei Punktalter lokalisieren keinen Gipfel.\par\smallskip
Hauptfolie 3: BM 50/51--75/76--100 gemeinsam mit je mindestens 500 Vertragsjahren.
Die günstigere BM-Mischung älterer Gruppen verdeckt dort einen Teil des Unterschieds.
\source{reports/s6a/existing\_age\_profiles.json; reports/s3/bm\_common\_cells.json, bm\_standardized.json, age\_contrasts.json (gepinnt).
Unterschiedliche Referenzpopulationen nicht direkt subtrahieren.}
\end{frame}

\begin{frame}[t]{B7 | Geografie: gleiche Zellen, verschiedene Populationen}
\begin{center}\small\StandardTable\end{center}
\smalltext Gemeinsame Zellen aus Alter, BM, Fahrzeugalter und Kraftstoff,
zusätzlich mit/ohne Dauer. Je geografischer Gruppe mindestens 100 Vertragsjahre und zehn Schadenzeilen je Zelle.\par\medskip
Ohne Dauer: 48 Zellen, 501\,612 Verträge. Mit Dauer: 79 Zellen, 389\,468 Verträge.
Gewichte: gepoolte Exponierungsanteile, je Vergleich für alle Gruppen gleich.\par\smallskip
EF--AB: Frequenzkontraste positiv; Schadenhöhen nicht eindeutig getrennt.
Verlustkontrast ohne Dauer: 0,68 [--146,67; 90,02] EUR/Jahr;
mit Dauer: 42,19 [7,78; 87,20].
\source{reports/s6a/standardized\_\{without,with\}\_duration.json, zugehörige Intervalle und support-Dateien.
500 bedingte Vertragsziehungen; Support und Gewichte fest. Keine isolierte Dauerinterpretation.
Region, Marke und Leistung bleiben in den Zellen verschieden.}
\end{frame}

\begin{frame}[t]{B8 | Datenherkunft und externer Methodenvergleich}
\smalltext
\textbf{Q1: CASdatasets, freMTPL2-Dokumentation}\par
Version 1.2-1, Commit 227fb56b8734; unbekannter Versicherer, ungefähr 2011--2013.
Area/Density: Wohnort; Region: historische Vertragsregion. Einige Zahlbeträge nach IRSA-IDA.
\href{https://github.com/dutangc/CASdatasets/blob/227fb56b8734bdb7c0327a41180e01d2ddaeaf26/man/freMTPL.Rd}{Originale Datensatzbeschreibung}\par\medskip
\textbf{Q2: ONISR, Bilan 2024; zeitnäherer Kontext: Bilan 2012}\par
Verkehrsteilnehmer aller Verkehrsarten, Frankreich métropolitaine; BAAC und
ONISR--Université Gustave Eiffel-Schätzungen für Verletzte.
Unfallort und körperliche Schwere, keine versicherten Wohnort-Zahlbeträge oder passenden Vertrags-/Kilometernenner.
2012 nur Übersichtsseite gelesen; widersprüchliche Gesamtzahlen nicht übernommen.
\href{https://www.onisr.securite-routiere.gouv.fr/etat-de-linsecurite-routiere/bilans-annuels-de-la-securite-routiere/bilan-2024-de-la-securite-routiere}{ONISR 2024}
\quad\href{https://www.onisr.securite-routiere.gouv.fr/etat-de-l-insecurite-routiere/bilans-annuels-de-la-securite-routiere/bilan-2012-de-la-securite-routiere}{ONISR 2012}\par\medskip
\textbf{Q3: Cerema, 20. Juli 2020}\par
80-km/h-Regel ab Juli 2018 auf Gegenverkehrsstraßen ohne Mitteltrennung;
Vorher/Nachher mit übrigem Netz als Kontrolle, Referenz 2013--2017,
Hauptnachperiode Juli 2018--Dezember 2019. Beobachtende Maßnahmenevaluation mit Vergleichsannahmen.
\href{https://www.cerema.fr/en/actualites/lowering-maximum-authorised-speed-80-kmh-final-assessment}{Originäre Ergebnis- und Methodendarstellung}
\source{Abruf der Originalseiten am 03.11.2026. Quellen und verwendete Abschnitte stehen in sources.json.
Kein verknüpfter Versicherungsnachweis. Nötig: Unfallort, Straßentyp, Kilometer, Fahrer, Personenschaden, Deckung und Regulierungsverlauf.}
\end{frame}

\begin{frame}[t]{B9 | Zwei Ursprungsmerkmale und ihre Zusatzinformation}
\begin{columns}[T,onlytextwidth]
\column{.54\textwidth}
{\smalltext\AllFeatureTable}\par\medskip
\smalltext Dummy-/Spline-Spalten gehören zum Ursprungsmerkmal.
Exponierung skaliert die Periode, ist kein drittes Tarifmerkmal.
\column{.43\textwidth}
\smalltext Ursprünglich vier Paare mit festem Alter;
danach Einzelmerkmale und BM + Area ergänzt.
Gleiche fünf Folds, jeweils Trainingsclaim-Mittel.\par\medskip
BM-GLM minus BM-GAM (je mit Alter):
0,593 [--0,188; 1,959]. Modellklassen nicht eindeutig getrennt.\par\medskip
Gewichtete Korrelation Alter--BM: \AgeBMRho;
Spearman ungewichtet: \AgeBMRank.\par\smallskip
Area-$\eta^2$: Alter \AreaAgeEta, BM \AreaBMEta.
Andere Skala; kleine Werte beweisen keine Unabhängigkeit oder stabileren Effekte.
\end{columns}
\source{reports/s6a/scores.json, bootstrap\_best.json; ablation/\{scores,dependence,report\}.json.
Explorativer Kandidatenkreis, keine globale Paaroptimalität. Fehlende Refit-/Concurvity-Prüfung betrifft Interpretierbarkeit und Effektstabilität.}
\end{frame}

\begin{frame}[t]{B10 | Was die Intervalle erfassen und was fehlt}
\smalltext
\textbf{Kernstudie und Merkmalsvergleich:}
500 gepaarte Ziehungen ganzer Verträge, einschließlich Nullschäden.
Gleiche gezogene Indizes für beide Regeln und gezogener Exponierungsnenner;
2,5-/97,5-Perzentile. Alle Vorhersagen bleiben fest.\par\medskip
\textbf{O2:} $z_i=d(N_i,e_i\widehat r_i^{\rm neu})-d(N_i,e_i\widehat r_i^{\rm ref})$,
$\widehat\Delta=\sum_i z_i/\sum_i e_i$.
\[
u_i=z_i-\widehat\Delta e_i,\quad
\widehat{\rm SE}(\widehat\Delta)=\frac{\sqrt{n s_u^2}}{\sum_i e_i},\quad
\widehat\Delta\pm1{,}96\,\widehat{\rm SE}
\]
Asymptotische Verhältnisstatistik bei geeigneten unabhängigen Beiträgen und endlichen zweiten Momenten;
kein Bootstrapintervall.\par\medskip
\textbf{Nicht enthalten:} erneutes Training, Modellauswahl, Generatorunsicherheit,
unbekannte Haushaltsabhängigkeit und bisher unbeobachtete Extremverluste.
Überlappende Trainingsfolds und drei Seeds sind keine unabhängigen Portfolios.
\source{reports/s4/paired\_bootstrap.json; reports/s5/evaluation\_bootstrap.json; reports/s6a/bootstrap\_best.json;
src/augmentation.py: paired\_difference; paper Anhang C. Bedingte Streuungsdiagnosen, keine umfassenden Tarifintervalle.}
\end{frame}

\begin{frame}[t]{B11 | CTGAN: Sampling und gemeinsame Unterstützung}
\smalltext
\textbf{Design:} CTGAN 0.11.1, CPU, 30 Epochen; nur Trainingsmerkmale,
Anzahl und positive Exposure. Keine IDs, Beträge oder Evaluationszeilen.
Zusatzratio zählt Zeilen, nicht Vertragsjahre.\par\medskip
\textbf{Sampling:} Seltene Kategorien werden im Training logarithmisch übergewichtet.
Der Adapter setzt im verwendeten Release die Samplingwahrscheinlichkeiten danach
auf empirische Trainingshäufigkeiten zurück. Das garantiert keine korrekte gemeinsame Verteilung.\par\medskip
\textbf{Ablehnung:} 102\,152 von 600\,000 Kandidaten verworfen (17,03\,\%);
Werte müssen gültig und im marginalen Trainingssupport liegen. Ablehnung verändert die Verteilung.\par\medskip
\textbf{Gemeinsamer Support:} \SupportMin--\SupportMax\,\% der akzeptierten Zeilen
liegen außerhalb des beobachteten Density-Bereichs ihrer Area im Trainingssubset.
Nachträgliche Diagnose, keine physische Unmöglichkeitsprüfung und keine neue Filter-/Fitregel.\par\medskip
Globale Gewichtsnormierung ändert die normierte GLM-Strafe nicht;
Bootstrap-Kontrollen messen stochastische Umgewichtung.
Zusätzliche Zeilen sind keine unabhängige Information und bieten keine Datenschutzgarantie.
\source{reports/o2/report.json, joint\_support.json, sources.json; config/augmentation\_protocol.json;
src/augmentation.py; Xu et al. (2019), [M4]. Keine Behauptung von Generator-Konvergenz oder Differential Privacy.}
\end{frame}

\begin{frame}[t]{B12 | Kosten, Vertrieb und Übertragung nach Deutschland}
\[
P_{\rm fix}=q+50,\qquad P_{\rm var}=q+0{,}15P_{\rm var}=q/0{,}85
\]
\smalltext Beide hypothetisch in EUR/Jahr. Schnittpunkt $q=850/3=283{,}33\ldots$ EUR/Jahr:
darunter variable Regel günstiger, darüber fixe Regel.
15\,\% Prämienanteil entspricht $1/0{,}85-1\approx17{,}65$\,\% Aufschlag auf $q$.\par\medskip
Mit definierten jährlichen Kosten $F$, Risiko-/Kapitallast $R$, Gewinnbetrag $G$
und variablem Prämienanteil $v$: $P=(q+F+R+G)/(1-v)$.
Regulierungskosten abgrenzen, um Doppelansätze zu vermeiden.\par\medskip
\textbf{Vertrieb:} Ein günstigerer Direktpreis kann bei tatsächlich niedrigeren Vollkosten denselben Deckungsbeitrag liefern.
Akquisitionskosten, Konversion und Selektion können das umkehren.
freMTPL2 identifiziert weder Kanalwirkung noch Profitabilität.\par\medskip
\textbf{Deutschland:} aktueller Zielbestand, gleiche Deckung/Haftung, Schadeninflation,
Bewertungspraxis, Nutzung und zeitliche Validierung fehlen.
Französischer Bonus-Malus ist keine direkt übertragbare SF-Klasse.
\source{reports/s6a/costs.json; docs/aktuarische\_einordnung.md; paper Anhang A.
Keine tatsächlichen Kosten, Gewinnspannen, Elastizitäten oder Steuersätze geschätzt. Die Gleichungen sind Szenarien.}
\end{frame}

\begin{frame}[t]{B13 | Methodische Literatur und nachprüfbare Ergebnisse}
\smalltext
\textbf{M1} scikit-learn: \href{https://scikit-learn.org/1.9/modules/linear_model.html}{GLM-Zielfunktion}
und \href{https://scikit-learn.org/1.9/modules/generated/sklearn.ensemble.HistGradientBoostingRegressor.html}{Histogramm-Boosting}.
Projektparameter in den registrierten JSON-Protokollen.\par\medskip
\textbf{M2} Eilers \& Marx (1996): \emph{Flexible Smoothing with B-splines and Penalties}.
\href{https://doi.org/10.1214/ss/1038425655}{Statistical Science 11(2), 89--102}.\par\smallskip
\textbf{M3} Wood: \href{https://stat.ethz.ch/R-manual/R-devel/library/mgcv/html/identifiability.html}{Identifiability constraints}.
mgcv als methodische Referenz, keine verwendete Fitbibliothek.\par\medskip
\textbf{M4} Xu et al. (2019): \emph{Modeling Tabular Data using Conditional GAN}.
\href{https://proceedings.neurips.cc/paper_files/paper/2019/file/254ed7d2de3b23ab10936522dd547b78-Paper.pdf}{NeurIPS 32}.\par\smallskip
\textbf{M5} Gneiting (2011): \emph{Making and Evaluating Point Forecasts}.
\href{https://doi.org/10.1198/jasa.2011.r10138}{JASA 106, 746--762}.\par\smallskip
\textbf{M6} Duan (1983): \emph{Smearing Estimate: A Nonparametric Retransformation Method}.
\href{https://doi.org/10.1080/01621459.1983.10478017}{JASA 78, 605--610}.\par\medskip
Zahlen, Tabellen und Diagrammwerte: \texttt{slides/evidence.json}.
Vollständige Literatur: \texttt{paper/literatur.bib}; Herleitungen: Begleitpaper.
Externe Quellen und tatsächlich gelesener Umfang: \texttt{slides/sources.json}.
\source{Ergebnisdateien: reports/s6/results.json, reports/s6a/ und reports/o2/.
Alle drei Vortrags-PDFs stammen aus derselben Beamer-Quelle. Daten und Modelle werden beim Folienbau nicht neu berechnet.}
\end{frame}

\fi
\end{document}
